% English translation of checked_pages/esq_119_035.tex; source: cote 119, site 35, PDF page 36, typescript page 33.
% AI-translated by GPT-6 Astra (OpenAI, ChatGPT Pro mode).
\scanpage{119/35}
\mspage{33}

over number fields, etc. The introduction of a subfield
$K \subset \mathbb{R}$ associated with the theory $\hand{\mathcal{T}}$ (formed by the points of $\mathbb{R}$ that are
$T$-\uline{tame}, i.e. such that the corresponding one-point set
is tame) makes it possible to introduce, for every point $x$ of a tame space
$X$, a residue field $k(x)$, which is a subextension of $\mathbb{R}/K$
algebraically closed in $\mathbb{R}$, and of finite transcendence degree
over $K$ (bounded above by the topological dimension of $X$). When the
transcendence degree of $\mathbb{R}$ over $K$ is infinite, one obtains a
notion of the transcendence degree (or ``dimension'') of a point of a
tame space, close to the familiar notion in algebraic
geometry. Such notions are absent from ``semianalytic'' tame
topology, which, on the other hand, appears to be the most appropriate
topological context for including real
and complex analytic spaces.

Among the first theorems one expects in
a tame topology as I envisage it, apart from comparison
theorems\frapparatus{V-FR-035-001}, are statements establishing, in an
appropriate sense, the existence and uniqueness of ``the'' tubular neighborhood
of a closed tame subspace in a tame space (compact, for
simplicity), the concrete ways of obtaining it (for example from
any tame map\frapparatus{V-FR-035-002} $X \longrightarrow \mathbb{R}^{+}$ having $Y$ as
its zero set), the description of its ``boundary'' (although in
general it is by no means a manifold with boundary !) $\partial T$, which has
in $T$ a neighborhood isomorphic to the product of $T$ with an interval,
etc. Under suitable equisingularity hypotheses, one
expects $T$ to be endowed, in an essentially unique way,
with the structure of a locally trivial bundle over $Y$, having $\partial T$
as a subbundle. This is one of the least clear points in
my provisional intuition of the situation, whereas the
homotopy class of the anticipated structural map $T \longrightarrow Y$
has an obvious meaning, independently of any equisingularity
hypothesis, as a homotopy inverse of the inclusion map $Y \longrightarrow T$,
which must be a homotopy equivalence. One way of obtaining
such a structure a posteriori would be through the hypothetical equivalence of
isotopy categories envisaged at the beginning, taking into account the fact that

\margin{35}
